\documentclass[12pt,a4paper]{article}

\usepackage[UTF8]{ctex}
\setCJKmainfont{Noto Serif CJK SC}
\usepackage{geometry}
\geometry{left=1.91cm,right=1.91cm,top=2.54cm,bottom=2.54cm}
\usepackage{amsmath}
\usepackage{booktabs}
\usepackage{hyperref}
\usepackage{float}

\title{\textbf{遣唐使交流网络的知识图谱分析}}
\author{}
\date{}

\begin{document}

\maketitle

\vspace{-1.5em}
\begin{center}
\small 分析代码详见 \url{https://github.com/szhhwh/kento-shi-network}
\end{center}

\section{研究问题}

遣唐使是公元7至9世纪日本向唐朝派遣的官方使团，持续了两个多世纪。除了外交任务，遣唐使也承担了大量文化学习的功能：留学僧和留学生随使团入唐，学习佛法、制度、技术，归国后将这些知识带回日本。

过去的研究大多聚焦于个人。空海在长安青龙寺学密、最澄在天台山求法、阿倍仲麻吕在唐为官与王维、李白等交游，这些故事已经有很详细的考证。但换个角度看：如果把两百多年间所有遣唐使的活动放在一起，能看到什么样的整体图景？具体而言，哪些人在知识传递中起到了连接不同群体的作用？遣唐使内部是否存在不同的交流群体？遣唐使的活动重心有没有随着时间发生变化？

为考察以上问题，使用知识图谱和社会网络分析的方法。数据来自七种文献：日本六国史\footnote{包括《日本書紀》《続日本紀》《日本後紀》《続日本後紀》《日本文徳天皇実録》《日本三代実録》。}约190卷编年体记录、中国正史中的新旧唐书约425卷、圆仁的《入唐求法巡礼行记》4卷日记、《善邻国宝记》3卷外交档案，以及空海的《御请来目录》、淡海三船的《唐大和上东征传》（鉴真传记）、和《入唐五家传》中的慧运传。空海目录详细列出了他从唐朝带回的密教经典和法具；鉴真传记录了鉴真六次东渡的过程；慧运传补充了入唐僧人的行程记录。这七种文献代表了日本官方的、中国官方的、个人旅行者的、求法僧自身的和外交档案的不同视角，综合起来才能看到相对完整的图景。

\section{分析方法}

分析过程分为三步。第一步是从文献中提取事件。先从全部文献文本中找出跟遣唐使活动相关的段落。检索时，使用面向古汉语预训练的语言模型\footnote{SIKU-BERT/sikuroberta，\url{https://huggingface.co/SIKU-BERT/sikuroberta}}将全部文本段和查询语句编码为向量，然后通过向量相似度找出与每句查询最相关的段落。查询的设计覆盖了遣唐使活动的五个主要维度（表\ref{tab:queries}）：派遣与出使、文化传入、师从学习、制度与技术传入、以及路线与地点。每个维度设计6至8句查询语句，用现代汉语描述并搭配“遣唐使”“入唐”“求法”“密教”等关键术语来引导模型定位相关段落。

\begin{table}[H]
\centering
\caption{五组语义查询的设计}
\label{tab:queries}
\begin{tabular}{p{2.5cm}p{3.5cm}p{6cm}}
\toprule
\textbf{查询组} & \textbf{目标} & \textbf{查询示例} \\
\midrule
派遣与出使 & 遣唐使的任命、出发、使团组成 & “遣唐使大使副使的任命与派遣”“日本遣唐使船从难波出发” \\
文化传入 & 佛教经典、佛像、制度等传入日本 & “遣唐使从唐朝带回佛经佛像”“最澄空海将密教传入日本” \\
师从学习 & 留学僧师从唐人、受法受戒 & “日本僧人在唐朝寺院学习佛法”“圆仁在五台山求法巡礼” \\
制度与技术 & 历法、官制、医学、建筑等传入 & “遣唐使引入唐朝律令制度”“日本模仿唐朝城市建设” \\
路线与地点 & 航行路线、登陆港口、活动地点 & “遣唐使到达长安洛阳”“圆仁在扬州的见闻” \\
\bottomrule
\end{tabular}
\end{table}

由于古文中“遣使”一词同样出现在遣新罗使、遣渤海使等不同外交活动的记录中，仅靠语义搜索可能引入不相关的内容，因此进一步用“遣唐”“入唐”“大唐”等遣唐使特有的关键词进行筛选。筛选后得到约1,100个相关段落。表\ref{tab:sources}列出了各文献的文本段分布情况。

\begin{table}[H]
\centering
\caption{各文献的文本段统计}
\label{tab:sources}
\begin{tabular}{lrr}
\toprule
\textbf{文献} & \textbf{文本段数} & \textbf{章节数} \\
\midrule
旧唐书·新唐书 & 25,310 & -- \\
六国史 & 8,293 & 190 \\
入唐求法巡礼行記 & 560 & 4 \\
善邻国宝记 & 195 & 3 \\
御请来目录 & 41 & -- \\
唐大和上东征传 & 61 & -- \\
入唐五家传·慧运 & 5 & -- \\
\midrule
\textbf{合计} & \textbf{34,465} & \textbf{197} \\
\bottomrule
\end{tabular}
\end{table}

然后从这些段落中逐一提取事件。由于筛选后的段落数量较多，约1,100段，采用分批并行处理的方式：首先为事件抽取任务编写详细的系统提示词，在其中逐一说明十四种事件类型的定义、关键词和正反示例，并明确规定提取原则：每条事件必须有直接的原文依据，模糊的记载，例如只说了年份但未提具体人名，不能猜测。然后将筛选出的段落分成若干批次，每批8个段落，每个批次包含系统提示和该批的全部文本，分别派发给多个 Agent 并行处理。每个Agent按照统一的提示词逐段分析，以结构化的格式输出每条事件的五元组信息：主体（谁）、动作（做了什么）、客体（对谁/对什么）、地点（在哪里）、时间（何时）。对于没有匹配任何事件类型的段落，要求 Agent 返回空结果。所有批次处理完成后，收集并合并各 Agent 的输出，按事件类型、主体、客体和原文前30字进行去重。最终得到742条事件，各类事件的分布如表\ref{tab:event_types}所示。

系统提示中为每种事件类型给出了明确的定义、关键词和正反例。以下是其中两种类型的示例：

\begin{quote}
\small
\textbf{1. 派遣遣唐使}\\
定义：官方任命大使/副使、遣唐使团正式出发的记载。必须有任命或使团出发的明确动作。\\
示例：天皇任命藤原常嗣为遣唐大使。\\
关键词：任命大使、拜遣唐使、为遣唐大使、遣唐使出发、节刀、遣于唐国\\
注意：如果原文只说某人“入唐”“赴唐”而没有任命或使团出发，应归类为“到达唐朝”或“前往地点”，而非此类型。\\[4pt]
\textbf{4. 师从学习}\\
定义：日本人在唐朝明确向某人学习的记载，要求原文有明确的师从或学习关系。\\
示例：空海师从惠果学习密教。\\
关键词：师从、从某学、受业、师事、禀受\\[4pt]
\textbf{抽取原则}：每条事件必须有原文依据；不可仅因词语共现就判定关系；使用原文中的名称，不翻译改写；模糊记载不猜测；没有事件则返回空。
\end{quote}

\begin{table}[H]
\centering
\caption{各类事件的分布}
\label{tab:event_types}
\begin{tabular}{lrlr}
\toprule
\textbf{事件类型} & \textbf{数量} & \textbf{事件类型} & \textbf{数量} \\
\midrule
到達唐朝 & 201 & 返回日本 & 98 \\
派遣遣唐使 & 93 & 接受賞賜 & 75 \\
海上遭難 & 90 & 授予官職 & 38 \\
前往地點 & 38 & 朝貢獻物 & 26 \\
受法受戒 & 26 & 帶回文化物品 & 26 \\
師從學習 & 13 & 傳入制度 & 13 \\
賦詩唱和 & 3 & 創立宗派 & 2 \\
\bottomrule
\end{tabular}
\end{table}

第二步是从事件中构建网络。构建的过程分两步。首先从每条事件中收集实体：事件的主体（谁）、客体（对谁/对什么）和地点（在哪里）都被当作独立的实体节点。实体的类型通过名称中的特征字来推断：例如包含“寺”“山”“州”“京”等字的判定为地点，包含“经”“律”“宗”“法”等字的判定为文化对象，包含“船”“舶”“团”等字的判定为群体。对于圆仁、空海、最澄等已知的历史人物，通过白名单确保不会被误判为其他类型。此外，对提取出的实体名进行了碎片化归并处理：例如将“最澄法師”“最澄賜”等变体统一归并为“最澄”，避免同一人物因称谓差异被当作多个不同实体。最终得到1,492个实体节点，分为人物、事件、地点、文化对象和群体五种类型。

然后，从每条事件中生成两类边。第一类是事件到实体的元数据边：从事件向主体连一条“事件主体”边、向客体连一条“事件客体”边、向地点连一条“事件地点”边，共1,394条，用于记录事件的时间与空间归属。第二类是将事件描述的关系直接转化为实体之间的社交边：事件类型决定了边类型——“前往地点”事件生成从人物到地点的“出行”边，“师从学习”或“受法受戒”事件生成从人物到人物或文化对象的“学习”边，“带回文化物品”等事件生成从人物到文化对象的“引入”边，“派遣遣唐使”或“朝贡献物”事件生成从人物到事件的“参与”边。这522条社交边反映人物与人物、人物与地点/文化对象之间的直接关系，是社会网络分析的核心数据。边分布如表\ref{tab:edges}所示。

\begin{table}[H]
\centering
\caption{知识图谱中的实体分布}
\label{tab:entities}
\begin{tabular}{lr}
\toprule
\textbf{实体类型} & \textbf{数量} \\
\midrule
Person（人物） & 319 \\
Event（事件） & 742 \\
Place（地点） & 319 \\
Culture（文化对象） & 59 \\
Group（群体/船只） & 53 \\
\midrule
\textbf{合计} & \textbf{1,492} \\
\bottomrule
\end{tabular}
\end{table}

\begin{table}[H]
\centering
\caption{知识图谱中的边分布}
\label{tab:edges}
\begin{tabular}{lrlr}
\toprule
\multicolumn{2}{c}{\textbf{社交边（522）}} & \multicolumn{2}{c}{\textbf{元数据边（1,394）}} \\
\midrule
出行 & 321 & 事件主体 & 657 \\
参与 & 120 & 事件地点 & 522 \\
学习 & 45 & 事件客体 & 215 \\
引入 & 36 & & \\
\midrule
\multicolumn{2}{l}{\textbf{合计}} & \multicolumn{2}{l}{\textbf{1,916}} \\
\bottomrule
\end{tabular}
\end{table}

这里有一个需要处理的问题。所有遣唐使都去过日本和唐朝。如果直接拿“去过日本”作为两个人之间的连接，那几乎所有人都会被连在一起，这种因为共享了通用地点而产生的连接其实没有分析价值。为了解决这个问题，在计算两个人之间的关联强度时，看的是他们共同去过的地方占各自去过的地方的比例：

\begin{equation}
J(A,B) = \frac{|P_A \cap P_B|}{|P_A \cup P_B|}
\label{eq:jaccard}
\end{equation}

其中 $P_A$ 和 $P_B$ 分别为人物 $A$ 和 $B$ 关联的地点集合或文化对象集合。按照公式(\ref{eq:jaccard})，如果两人仅因到过日本而产生交集，但各自还去过多个其他地点，分母中的并集因这些差异而增大，权重自然降低；反之，若两人均前往五台山、长安等多处特定地点，则权重较高，反映更可能的历史关联。这一方法避免了仅由日本、唐朝等通用地点造成的虚假连接。

第三步是在这个加权网络上做分析，计算了以下四个指标。

\textbf{中介中心性}衡量一个节点在网络中作为“桥梁”的程度。对于节点 $v$，其定义为经过 $v$ 的任意两节点间最短路径数占所有最短路径数的比例：

\begin{equation}
C_B(v) = \sum_{s \neq v \neq t} \frac{\sigma_{st}(v)}{\sigma_{st}}
\label{eq:betweenness}
\end{equation}

其中 $\sigma_{st}$ 为节点 $s$ 与 $t$ 之间最短路径的总数，$\sigma_{st}(v)$ 为其中经过 $v$ 的路径数。中介中心性高的人，意味着网络中有较多的信息流动需要经过他；在此处的历史语境中，反映的是该人物连接了较多本不直接相关的其他人物。

\textbf{加权度}衡量一个节点关联的边的权重之和：

\begin{equation}
D_w(v) = \sum_{u \in N(v)} w_{uv}
\label{eq:degree}
\end{equation}

其中 $N(v)$ 为 $v$ 的邻居集合，$w_{uv}$ 为 Jaccard 加权后的边权重。加权度高的人，意味着他关联的地点或文化对象较为多样且独特。

\textbf{社区检测}使用 Louvain 算法，通过最大化模块度 $Q$ 来发现网络中的自然聚类：

\begin{equation}
Q = \frac{1}{2m} \sum_{ij} \left[ w_{ij} - \frac{k_i k_j}{2m} \right] \delta(c_i, c_j)
\label{eq:modularity}
\end{equation}

其中 $w_{ij}$ 为边权重，$k_i$ 为节点 $i$ 的加权度，$m$ 为总边权重，$\delta(c_i, c_j)$ 在 $i$ 与 $j$ 同属一个社区时为1、否则为0。$Q$ 越接近1，说明网络的社区结构越显著。

\textbf{学习传播比}（L/I 比）定义为学习边数与引入边数之比，用于衡量网络的功能侧重：

\begin{equation}
R_{L/I} = \frac{|\{\text{学习边}\}|}{|\{\text{引入边}\}|}
\label{eq:li}
\end{equation}

当 $R_{L/I} > 1$ 时，网络中学习关系多于传播关系，表明以知识获取为主要功能；当 $R_{L/I} < 1$ 时，则表明文化传播更为活跃。

最后，为了看时间上的变化，需要将742条事件标注到遣唐使的各批次上。每批遣唐使有记载的出发和归国年份，以及大使、副使等人员名单。标注采用多轮匹配策略。

第一轮是直接匹配：从每条事件的原文和时间信息中查找“承和”“天平”“延历”等日本年号，根据“年号+年份”换算出公元年份，再匹配到在时间上最接近的遣唐使批次。例如“承和五年”换算为公元838年，匹配到838年出发的遣唐使。同时，如果事件的原文中出现了某批遣唐使的大使或副使姓名，则直接标注到对应批次。

第二轮是上下文回溯。由于第一步中将全文切为200字一段，年号可能出现在相邻的段落中而与事件正文分离。例如某段只有“到登州入开元寺宿”，但前一段中有“开成五年”。因此对第一轮未匹配的事件，检查其前后各8段的文本，寻找丢失的年号信息。

第三轮利用人物的已知活动时间来推断。对于前两轮仍未匹配的事件，如果事件的主体在笔者整理的遣唐使人物生平数据库中有记录——该数据库从日本六国史、《入唐求法巡礼行记》等原始史料及相关研究专著中收录了78位有明确生卒或入唐、归国年份记载的人物——则根据其入唐和归国的年份来推断所属批次。经过三轮匹配，742条事件中约87\%完成了时间标注。

\section{网络中的核心人物与结构}

\subsection{文献记录中的广泛旅行者}

中介中心性最高的是圆仁，其次是空海（表\ref{tab:centrality}）。圆仁的《入唐求法巡礼行记》近乎逐日记载了他从承和五年（838年）入唐到承和十四年（847年）归国的全部行程：到扬州海陵县、登州赤山浦、五台山、长安大兴善寺和青龙寺，以及沿途几十个地点。

\begin{table}[H]
\centering
\caption{中介中心性最高的十位人物}
\label{tab:centrality}
\begin{tabular}{clcc}
\toprule
\textbf{排名} & \textbf{人物} & \textbf{中介中心性} & \textbf{加权度} \\
\midrule
1 & 円仁 & 0.00207 & — \\
2 & 入唐學問僧 & 0.00065 & — \\
3 & 空海 & 0.00065 & — \\
4 & 智通·智達二師 & 0.00005 & — \\
5 & 釋道昭·道嚴等十三人 & 0.00005 & — \\
6 & 僧永忠 & — & 4.00 \\
7 & 円載留學傔從僧仁濟 & — & 4.00 \\
8 & 大僧都行賀 & — & 4.00 \\
9 & 平郡朝臣廣成 & — & 4.00 \\
10 & 釋榮睿·普照 & — & 4.00 \\
\bottomrule
\end{tabular}
\end{table}

空海排名第三同样有充分的史料基础。空海的《御请来目录》详细记录了他在长安青龙寺师从惠果阿阇梨的过程，以及受法、灌顶的具体细节。空海在目录中写道自己“沐受明灌顶再三焉，受阿阇梨位一度也”。这些自体记录在事件抽取中产生了多条高质量的师从学习和受法受戒事件。

值得注意的是，《唐大和上东征传》（鉴真传）和《入唐五家传·慧运传》虽然提供了丰富的传记材料，但由于鉴真主要是渡日传法而非入唐求法，慧运的路线也相对独立，他们在人物-地点网络中的中介性不如圆仁和空海突出。这反映了分析方法本身的特性：基于地点共现的Jaccard投影天然对行程广泛、停留地点多样的人物更敏感。

\subsection{遣唐使是一批一批组织的}

人与人之间的直接连接几乎全在同一批遣唐使的内部。天平胜宝四年（752年）那批遣唐使，以大使藤原清河、副使大伴古麻吕和吉备真备为首，之间有很多共现，但他们跟承和五年（838年）入唐的圆仁、藤原常嗣这批人几乎没有直接联系。从Jaccard加权后的网络来看，不同批次的遣唐使成员虽然可能去过相同的地方，但在人际层面高度隔离，仅通过阿倍仲麻吕等少数长期留唐者存在有限的跨批次联系。

这其实很好理解。遣唐使的航海风险很大：日本史料中频繁出现“风波急激，船被打破”“船触礁覆没”的记载。每次派遣的规模和时间都有明确限制，通常两到四艘船，在唐逗留数月至数年，每一批实际上就是一个独立的团队，跟前后批次之间缺少人际延续。遣唐使的网络结构是由其制度安排决定的：各批次独立成行、独立返回，而非一个连续运转的社交网络。

事件的时间分布进一步支持了这一点。末期集中了413条事件。838年，藤原常嗣率最后一次实际成行的遣唐使入唐，圆仁随行并在唐九年间留下了逐日记录，贡献了该时期最多的事件。这种集中在时间上与遣唐使制度的组织方式是一致的。

\subsection{地点网络的扩展}

空海目录、鉴真传和慧运传提供了丰富的地点信息。空海目录中记录了长安西明寺、青龙寺等寺院；鉴真传中记录了扬州大明寺、明州阿育王寺、台州国清寺等地点；慧运传中记录了肥前国松浦郡、大唐温州乐城县等港口。这些地点丰富了地点网络中关于寺院和登陆港口的信息。

从时间上看，地点网络的扩展呈现出明显的阶段性：初期地点以日本出发港大宰府、博多和唐朝主要城市长安、洛阳为主，带有政治外交色彩；盛期开始出现更多的登陆地点和沿途寺院；到了末期，五台山、青龙寺、大兴善寺等佛教圣地成为网络中最活跃的节点。这个变化与遣唐使从以外交学习为主到以佛教求法为主的整体趋势是吻合的。

\section{遣唐使功能在时间上的变化}

\subsection{遣唐使始终以学习为核心功能}

把遣唐使按时期分组后（表\ref{tab:temporal}），比较每个时期的学习关系数和传播关系数之比。全部四个时期的L/I比均大于或等于1.0，表明遣唐使在任何时期都以学习为核心功能，这与遣唐使制度设立的初衷，吸收唐朝文化是一致的。初期L/I比为1.00，学习和传播基本平衡，反映早期遣唐使在派遣学问僧求法的同时，也带回了历法、制度和典籍。盛期L/I比升至1.47：吉备真备归国后“献《唐礼》一百三十卷”（《续日本纪》天平七年四月辛亥条），体现了既有学习又有制度性引入的特征；而阿倍仲麻吕终生仕唐未归，代表了另一种遣唐使成员的人生轨迹。中期L/I比达到最高的2.00，这与空海、最澄等精英僧侣的入唐求法高度吻合：他们回国后创立的日本天台宗和真言宗，虽然表现为“传播”性质的创立宗派事件，但其基础是在唐朝的数年密集学习。晚期L/I比为1.48，学习仍占主导，但传播活动，如带回物品、传入制度等，也在同步增长，反映出末期遣唐使在求法的同时大量搜集经论文书的特征。

\begin{table}[H]
\centering
\caption{四个时期的比较}
\label{tab:temporal}
\begin{tabular}{lcccccp{4cm}}
\toprule
\textbf{时期} & \textbf{年份} & \textbf{事件数} & \textbf{学习/传播比} & \textbf{网络密度} & \textbf{特点} \\
\midrule
初期 & 630--701 & 69 & 1.00 & 0.054 & 政治外交为主，学习与传播平衡 \\
盛期 & 702--770 & 146 & 1.47 & 0.052 & 学习活动显著增多 \\
中期 & 771--834 & 24 & 2.00 & 0.087 & 精英求法，学习高度集中 \\
晚期 & 835--894 & 413 & 1.48 & 0.030 & 事件量暴增，学习仍占主导 \\
\bottomrule
\end{tabular}
\end{table}

总体上，除了初期学习与传播基本持平外，盛期之后L/I比始终高于1.4，说明学习活动在遣唐使中占据稳定的主导地位。

\subsection{公元804年前后的变化}

分批次看数据，公元804年（日本年号延历二十三年）是一个重要的节点。空海和最澄同时入唐无疑是遣唐使历史上的标志性事件：空海在长安青龙寺师从惠果，最澄在天台山求法。从L/I比的时间序列看，学习活动的增长在盛期已经开始，中期达到顶峰，804年处于这一上升趋势之中。空海的《御请来目录》为此提供了第一手的史料支持：空海本人在目录中详细列举了从唐朝带回的“大日经”“金刚顶经”等密教经典和诸尊曼荼罗，光是经论章疏就有数百卷之多。鉴真传中记录的六次东渡过程，也为理解唐代中日佛教交流的广度和深度提供了珍贵的背景。

从事件类型的时间分布来看，受法受戒和传入制度两类事件在盛期就已经出现，但确实在末期显著增多。空海和最澄的入唐是这一长时段趋势中的高光时刻。

\section{局限性}

这个分析有三个比较明显的局限。一是圆仁的日记贡献了最多的事件，共42条源文本，远超其他人物，个人旅行日记的事件密度远高于官方编年史，体现了任何基于文本的分析都不可避免地受到文献存留偏差的影响：存世文献越丰富的人物，在量化分析中自然会占据更显眼的位置。

二是从古文中自动提取事件不可避免地存在误差。同名异人、异体字、复合地名的处理虽然经过了人工校核，但无法做到完全准确。分析中引入了严格的实体清洗规则，即过滤包含句读符号的文本碎片和过长的非人名实体，一定程度上提升了实体质量，但可能也过滤掉了一些合法的复合实体。

三是空海目录和鉴真传虽然丰富了史料基础，但这类私人记录、僧侣传记与六国史这类官方编年史有本质区别：前者细节丰富但覆盖面窄，后者覆盖面广但记录简略。两种性质不同的史料混合分析时，分析结果不可避免地会偏向细节丰富的文献。

\section{总结}

从以上分析中可以看到四点。第一，圆仁在文献记录中地位突出：他拥有最详尽的一手旅行记录，因此在知识图谱中占据中枢位置。同时，空海目录等文献也丰富了网络中其他核心人物的信息。这提示了一个方法层面的认知：基于文本的网络分析揭示的是“记录中可见的连接”而非“历史上存在过的所有连接”，对于文献存留极不均衡的古代史料，前者永远是后者的一个偏样本。第二，遣唐使是以使命批次为基本单位来组织的：每次派遣都是一个相对独立的小组，跟前后批次之间人际联系很少，仅有少数长期留唐者，如阿倍仲麻吕，起到有限的跨批次连接作用。这与当时的航海条件和使团规模吻合，表明网络结构不是偶然的，而是制度约束的自然结果。第三，多源文献的混合分析也带来了新的挑战：不同性质文献，如个人日记、官方编年史、僧侣目录，记录密度差异极大，不能简单等权重处理。第四，遣唐使始终以学习为核心功能：四个时期的L/I比均不低于1.0，学习集中度在空海、最澄入唐的中期达到最高。遣唐使功能的演变，体现在学习方式和内容的变化上：从早期政治制度与佛法并学，到盛期学问僧和留学生大批派遣，再到后期密教和天台宗的系统学习与宗派创立。

用知识图谱的方法来看遣唐使，好处是可以看到整体结构和长期趋势，而不只是一个人的故事。当然，这只是一种观察方式，不能替代对具体人物和事件的深入研究。两种方法结合起来，从宏观看到结构、从微观看清细节，可能才是理解这段历史最好的方式。

\end{document}
